% ============================================================================
% LANGENTWURF LE-01d: ¿Qué tal? + Stundenleistung
% Spanisch Klasse 7 - Sequenz 1: Empezamos
% ============================================================================
% Tier: 1 (vollständig)
% Doppelstunde: 1d (Std. 7-8 von 8) - ABSCHLUSS
% Fachfarbe: #c41e3a (Karminrot)
% ============================================================================

\documentclass[11pt,a4paper]{article}

\usepackage[utf8]{inputenc}
\usepackage[T1]{fontenc}
\usepackage[ngerman]{babel}
\usepackage{geometry}
\usepackage{fancyhdr}
\usepackage{lastpage}
\usepackage{xcolor}
\usepackage{tcolorbox}
\usepackage{enumitem}
\usepackage{tabularx}
\usepackage{longtable}
\usepackage{array}
\usepackage{booktabs}
\usepackage{amssymb}

% ============================================================================
% KONFIGURATION
% ============================================================================

\newcommand{\fach}{Spanisch}
\newcommand{\fachfarbe}{c41e3a}
\newcommand{\jahrgangsstufe}{7}
\newcommand{\schulart}{Gesamtschule}
\newcommand{\stundenthema}{¿Qué tal? + Stundenleistung}
\newcommand{\reihenthema}{Empezamos}
\newcommand{\stundennummer}{1d}
\newcommand{\reihenstunden}{4}
\newcommand{\gession}{A1}
\newcommand{\lehrwerk}{Qué pasa 1}
\newcommand{\lehrwerkseite}{S. 22-27}
\newcommand{\lerngruppengroesse}{24}

% ============================================================================
% LAYOUT
% ============================================================================
\geometry{left=2.5cm, right=2.5cm, top=2.5cm, bottom=2.5cm}
\definecolor{fach}{HTML}{\fachfarbe}
\definecolor{gray}{HTML}{666666}
\definecolor{lightgray}{HTML}{f5f5f5}
\definecolor{successgreen}{HTML}{2a7a4b}
\definecolor{warningorange}{HTML}{b35c00}
\definecolor{basis}{HTML}{2a7a4b}
\definecolor{standard}{HTML}{2c5aa0}
\definecolor{erweiterung}{HTML}{7b1fa2}

\tcbuselibrary{skins,breakable}

\newtcolorbox{infobox}[1][]{
    colback=lightgray,
    colframe=fach,
    fonttitle=\bfseries,
    title=#1,
    breakable
}

\newtcolorbox{scorebox}{
    colback=white,
    colframe=fach,
    fonttitle=\bfseries,
    title=Didaktik-Score,
    breakable
}

\pagestyle{fancy}
\fancyhf{}
\fancyhead[L]{\small\textcolor{gray}{\fach{} -- Jg. \jahrgangsstufe{} -- \stundenthema}}
\fancyhead[R]{\small\textcolor{gray}{LE-\stundennummer{} | Tier 1}}
\fancyfoot[C]{\small\thepage{} / \pageref{LastPage}}
\renewcommand{\headrulewidth}{0.4pt}

% Differenzierungsmarker
\newcommand{\B}{\textcolor{basis}{\textbf{[B]}}}
\newcommand{\Std}{\textcolor{standard}{\textbf{[S]}}}
\newcommand{\E}{\textcolor{erweiterung}{\textbf{[E]}}}

% ============================================================================
% DOKUMENT
% ============================================================================
\begin{document}

% ============================================================================
% DECKBLATT
% ============================================================================
\begin{titlepage}
    \centering
    \vspace*{2cm}
    {\large\textcolor{gray}{\schulart}}\\[2cm]
    {\Huge\textcolor{fach}{\textbf{Langentwurf}}}\\[0.5cm]
    {\large LE-\stundennummer{} | Tier 1 (vollständig)}\\[2cm]
    {\LARGE\textbf{\stundenthema}}\\[0.5cm]
    {\large Sequenz: \reihenthema{} -- \textbf{ABSCHLUSS}}\\[2cm]
    \begin{tabular}{rl}
        \textbf{Fach:} & \fach{} (A1) \\[0.3cm]
        \textbf{Klasse:} & \jahrgangsstufe \\[0.3cm]
        \textbf{Lehrwerk:} & \lehrwerk, \lehrwerkseite \\[0.3cm]
        \textbf{Doppelstunde:} & \stundennummer{} (Std. 7-8 von 8) \\[0.3cm]
    \end{tabular}
    \vfill
\end{titlepage}

\tableofcontents
\newpage

% ============================================================================
% 1. BEDINGUNGSANALYSE
% ============================================================================
\section{Bedingungsanalyse}

\subsection{Lerngruppe}
\begin{tabular}{ll}
    Kursgröße: & \lerngruppengroesse{} SuS \\
    Jahrgangsstufe: & \jahrgangsstufe \\
    Sprachniveau: & A1 (GER) -- Anfänger \\
    Fremdsprache: & 2. Fremdsprache (nach Englisch) \\
\end{tabular}

\subsection{Vorwissen aus 01a/01b/01c}
\begin{itemize}
    \item \textbf{01a:} Begrüßungen (Hola, Buenos días, Adiós)
    \item \textbf{01b:} Sich vorstellen (Me llamo..., Soy de..., ser-Konjugation)
    \item \textbf{01c:} Artikel (el/la/los/las), Alphabet, Buchstabieren
\end{itemize}

\subsection{Differenzierung (intern)}
\begin{tabular}{|l|p{3.5cm}|p{3.5cm}|p{3.5cm}|}
\hline
& \B{} \textbf{Basis} & \Std{} \textbf{Standard} & \E{} \textbf{Erweiterung} \\
\hline
\textbf{Ziel} & BR: Rezeption, Grundwortschatz & MR: Produktion mit Hilfe & GY: Freie Anwendung \\
\hline
\textbf{¿Qué tal?} & Antworten zuordnen & Fragen und antworten & + Nachfragen stellen \\
\hline
\textbf{Test} & Basisaufgaben (60\%) & Alle Pflichtaufgaben & + Zusatzaufgaben \\
\hline
\end{tabular}

% ============================================================================
% 2. SACHANALYSE
% ============================================================================
\section{Sachanalyse}

\subsection{Sprachliche Strukturen}

\textbf{Frage nach dem Befinden:}
\begin{center}
\begin{tabular}{|l|l|}
\hline
\textbf{Frage} & \textbf{Bedeutung} \\
\hline
¿Qué tal? & Wie geht's? (informell) \\
¿Cómo estás? & Wie geht es dir? (informell) \\
¿Cómo está usted? & Wie geht es Ihnen? (formell) \\
\hline
\end{tabular}
\end{center}

\textbf{Antworten:}
\begin{center}
\begin{tabular}{|l|l|l|}
\hline
\textbf{Spanisch} & \textbf{Deutsch} & \textbf{Stimmung} \\
\hline
Muy bien, gracias. & Sehr gut, danke. & ++ \\
Bien, gracias. & Gut, danke. & + \\
Regular. & Es geht so. & 0 \\
Así así. & So lala. & 0 \\
Mal. & Schlecht. & -- \\
\hline
\end{tabular}
\end{center}

\textbf{Rückfrage:}
\begin{itemize}
    \item ¿Y tú? = Und du? (informell)
    \item ¿Y usted? = Und Sie? (formell)
\end{itemize}

\subsection{Konsolidierung: Gesamtübersicht Sequenz 1}
\begin{itemize}
    \item \textbf{Begrüßung:} Hola, Buenos días/tardes/noches, Adiós, Hasta luego
    \item \textbf{Vorstellen:} Me llamo..., Soy de..., Mucho gusto
    \item \textbf{ser:} yo soy, tú eres, él/ella es, nosotros somos
    \item \textbf{Artikel:} el, la, los, las + Genus-Regeln
    \item \textbf{Alphabet:} 27 Buchstaben, besondere Aussprache
    \item \textbf{Befinden:} ¿Qué tal?, Bien/Mal/Regular, ¿Y tú?
\end{itemize}

% ============================================================================
% 3. DIDAKTISCHE ANALYSE
% ============================================================================
\section{Didaktische Analyse}

\subsection{Stellung in der Sequenz}
Dies ist die \textbf{4. und letzte Doppelstunde} (Std. 7-8 von 8) der Sequenz ``\reihenthema''.
Die Stunde schließt die Einführungssequenz ab mit:
\begin{enumerate}
    \item Neuem Inhalt: ¿Qué tal? (Befinden erfragen)
    \item Konsolidierung aller bisherigen Inhalte
    \item Stundenleistung (summative Bewertung)
\end{enumerate}

\subsection{Kompetenzziele (Rahmenplan MV)}
\begin{itemize}
    \item \textbf{Sprechen:} Einfache Gespräche über Befinden führen
    \item \textbf{Verfügen über sprachliche Mittel:} Grundwortschatz Begrüßung/Vorstellung
    \item \textbf{Interkulturell:} Höflichkeitsformen (tú vs. usted)
\end{itemize}

\subsection{Legitimation}
\textbf{Rahmenplan Spanisch MV:}
\begin{itemize}
    \item An Gesprächen teilnehmen: einfache Kontaktgespräche führen
    \item Verfügen über sprachliche Mittel: Redemittel für Alltagssituationen
\end{itemize}

% ============================================================================
% 4. STUNDENZIELE
% ============================================================================
\section{Stundenziele}

\subsection{Grobziel}
Die SuS erweitern ihre \textbf{kommunikative Kompetenz}, indem sie nach dem Befinden fragen und antworten können. Sie \textbf{konsolidieren} alle Inhalte der Sequenz 1 und weisen ihre Kompetenzen in einer \textbf{Stundenleistung} nach.

\subsection{Feinziele (operationalisiert)}
\begin{tabular}{|c|p{8cm}|c|c|}
\hline
\textbf{Nr.} & \textbf{Die SuS können...} & \textbf{AFB} & \textbf{Diff.} \\
\hline
FZ1 & mit ``¿Qué tal?'' nach dem Befinden fragen und mit passenden Redemitteln antworten & I/II & alle \\
\hline
FZ2 & eine kurze Begrüßungssequenz (Begrüßung, Vorstellen, Befinden, Verabschiedung) dialogisch durchführen & II & alle \\
\hline
FZ3 & die Inhalte der Sequenz 1 (Begrüßung, ser, Artikel, Alphabet) in einer Stundenleistung anwenden & II/III & alle \\
\hline
FZ4 & zwischen formellen (usted) und informellen (tú) Anredeformen unterscheiden & II & \Std{}/\E{} \\
\hline
\end{tabular}

% ============================================================================
% 5. METHODISCHE ANALYSE
% ============================================================================
\section{Methodische Analyse}

\subsection{Medien}
\begin{itemize}
    \item Tafel/Whiteboard
    \item Lehrwerk: \lehrwerk, \lehrwerkseite
    \item Arbeitsblatt: AB-01d.pdf (Übungen)
    \item Stundenleistung: SL-01d.pdf (Test)
    \item Audio: Track 10-11 (Dialoge)
    \item Optional: LP-01d.html (Wiederholung)
\end{itemize}

\subsection{Sozialformen}
\begin{itemize}
    \item Plenum: Einführung ¿Qué tal?, Wiederholung
    \item Partnerarbeit: Dialoge üben
    \item Einzelarbeit: Stundenleistung
\end{itemize}

\subsection{Besonderheit: Stundenleistung}
\begin{itemize}
    \item Dauer: 25-30 Minuten
    \item Umfang: 20 Punkte (4 Aufgaben)
    \item Bewertung: Note nach Prozentskala
    \item Differenzierung: Basis 60\% = ausreichend
\end{itemize}

% ============================================================================
% 6. VERLAUFSPLANUNG
% ============================================================================
\section{Verlaufsplanung}

\textbf{1. Stunde (45 Min.) -- Schwerpunkt: ¿Qué tal? + Übung}

\begin{longtable}{|p{1cm}|p{1.5cm}|p{4cm}|p{3cm}|p{2cm}|p{2cm}|}
\hline
\textbf{Zeit} & \textbf{Phase} & \textbf{Lehrerhandeln} & \textbf{Schülerhandeln} & \textbf{SF} & \textbf{Medien} \\
\hline
\endhead

0--5 & Einstieg & Begrüßung: ``¡Hola! ¿Qué tal?'' & Reagieren (vermuten) & Plenum & -- \\
\hline

5--15 & Input & ¿Qué tal? einführen, Antworten (bien/mal/regular) & Nachsprechen, Notieren & Plenum & Tafel \\
\hline

15--25 & Übung 1 & AB-01d Aufg. 1 (Zuordnung Antworten) & Zuordnen & EA & AB \\
\hline

25--35 & Übung 2 & Dialog-Kette: ¿Qué tal? -- Bien, ¿y tú? & Reihum antworten & Plenum & -- \\
\hline

35--45 & Sicherung & Merkkasten gemeinsam, HA erklären & Übertragen & EA/Plenum & AB \\
\hline
\end{longtable}

\vspace{1em}

\textbf{2. Stunde (45 Min.) -- Schwerpunkt: Wiederholung + Stundenleistung}

\begin{longtable}{|p{1cm}|p{1.5cm}|p{4cm}|p{3cm}|p{2cm}|p{2cm}|}
\hline
\textbf{Zeit} & \textbf{Phase} & \textbf{Lehrerhandeln} & \textbf{Schülerhandeln} & \textbf{SF} & \textbf{Medien} \\
\hline
\endhead

0--5 & Aktivierung & Schnelle Wiederholung: Artikel, ser & Antworten & Plenum & -- \\
\hline

5--10 & Vorbereitung & Stundenleistung erklären, Fragen klären & Zuhören, Fragen & Plenum & SL \\
\hline

10--35 & Test & Aufsicht, individuelle Hilfe bei Verständnisfragen & Stundenleistung bearbeiten & EA & SL \\
\hline

35--40 & Abgabe & Einsammeln, kurze Entspannung & Abgeben & -- & -- \\
\hline

40--45 & Abschluss & Ausblick Sequenz 2, Feedback & Rückmeldung & Plenum & -- \\
\hline
\end{longtable}

\textbf{Legende:} EA = Einzelarbeit, PA = Partnerarbeit, SF = Sozialform, SL = Stundenleistung

% ============================================================================
% 7. ERWARTETE SCHWIERIGKEITEN
% ============================================================================
\section{Erwartete Schwierigkeiten}

\begin{tabular}{|p{5cm}|p{8cm}|}
\hline
\textbf{Schwierigkeit} & \textbf{Reaktion} \\
\hline
Verwechslung ¿Qué tal? / ¿Cómo te llamas? & Bedeutungsunterschied klären: tal = Befinden, llamas = Name \\
\hline
Prüfungsangst bei Stundenleistung & Ruhige Atmosphäre, klare Instruktionen, Zeitzugabe für \B{} \\
\hline
Zeitdruck bei SL & Klare Zeitansagen, einfachere Aufgaben zuerst \\
\hline
ser-Konjugation vergessen & Kurze Wiederholung vor dem Test \\
\hline
\end{tabular}

% ============================================================================
% 8. HAUSAUFGABE
% ============================================================================
\section{Hausaufgabe}

\begin{itemize}
    \item \textbf{Vor der Stunde:} LP-01d.html zur Wiederholung (optional)
    \item \textbf{Nach der Stunde:} Vokabeln Sequenz 1 wiederholen für Vokabeltest
\end{itemize}

% ============================================================================
% 9. LÖSUNGEN
% ============================================================================
\section{Lösungen AB-01d}

\textbf{Merkkasten: ¿Qué tal?}
\begin{itemize}
    \item ¿Qué tal? = Wie geht's?
    \item Muy bien = Sehr gut
    \item Bien = Gut
    \item Regular / Así así = Es geht so
    \item Mal = Schlecht
    \item ¿Y tú? = Und du?
\end{itemize}

\textbf{Aufgabe 1 (Zuordnung):} (4P)
\begin{itemize}
    \item ++ $\rightarrow$ Muy bien, gracias
    \item + $\rightarrow$ Bien, gracias
    \item 0 $\rightarrow$ Regular / Así así
    \item -- $\rightarrow$ Mal
\end{itemize}

\textbf{Aufgabe 2 (Dialog vervollständigen):} (6P)
\begin{itemize}
    \item A: ¡Hola! ¿Qué tal?
    \item B: Bien, gracias. ¿Y tú?
    \item A: Muy bien. ¿Cómo te llamas?
    \item B: Me llamo [Name]. Mucho gusto.
\end{itemize}

\textbf{Aufgabe 3 (Wiederholung ser):} (4P)
\begin{itemize}
    \item yo soy
    \item tú eres
    \item él/ella es
    \item nosotros somos
\end{itemize}

\textbf{Aufgabe 4 (Artikel):} (6P)
\begin{itemize}
    \item el libro, la mesa, los chicos, las chicas
\end{itemize}

% ============================================================================
% 10. STUNDENLEISTUNG - ÜBERSICHT
% ============================================================================
\section{Stundenleistung SL-01d}

\textbf{Aufbau:} 4 Aufgaben, 20 Punkte, 25 Minuten

\begin{tabular}{|c|l|c|c|c|}
\hline
\textbf{Aufg.} & \textbf{Inhalt} & \textbf{AFB} & \textbf{Punkte} & \textbf{Diff.} \\
\hline
1 & Zuordnung: Begrüßungen & I & 4 & alle \\
\hline
2 & Lückentext: ser-Konjugation & I/II & 4 & alle \\
\hline
3 & Artikel einsetzen (el/la/los/las) & I/II & 6 & alle \\
\hline
4 & Dialog vervollständigen & II/III & 6 & alle \\
\hline
\multicolumn{3}{|r|}{\textbf{Gesamt:}} & \textbf{20} & \\
\hline
\end{tabular}

\textbf{Bewertung (Klasse 7):}
\begin{tabular}{|c|c|c|c|c|c|c|}
\hline
Note & 1 & 2 & 3 & 4 & 5 & 6 \\
\hline
ab \% & 96 & 80 & 60 & 40 & 20 & <20 \\
\hline
ab Punkte & 19 & 16 & 12 & 8 & 4 & <4 \\
\hline
\end{tabular}

% ============================================================================
% 11. DIDAKTIK-SCORE
% ============================================================================
\newpage
\section{Didaktik-Score}

\begin{scorebox}
\textbf{Bewertung der didaktischen Qualität nach 10 Dimensionen}\\
\textbf{Ziel-Score: $\geq$ 9.0 (Premium-Qualität)}

\vspace{1em}
\begin{tabular}{|c|l|c|c|p{4.5cm}|}
\hline
\textbf{\#} & \textbf{Dimension} & \textbf{Gew.} & \textbf{Score} & \textbf{Notiz} \\
\hline
1 & Differenzierung & 15\% & 9/10 & SL mit Basis-Sicherung, Zusatzaufgaben \\
\hline
2 & Taxonomie (AFB) & 10\% & 9/10 & AFB I: 40\%, AFB II: 45\%, AFB III: 15\% \\
\hline
3 & Lernziel-Alignment & 10\% & 10/10 & Rahmenplan: Kontaktgespräche \\
\hline
4 & Aktivierung & 10\% & 9/10 & Dialog-Kette, PA, SL \\
\hline
5 & Scaffolding & 10\% & 9/10 & Input $\rightarrow$ Übung $\rightarrow$ Test \\
\hline
6 & Feedback-Qualität & 10\% & 9/10 & Bewertungsraster transparent \\
\hline
7 & Sprachliche Klarheit & 10\% & 10/10 & Altersgerecht Kl. 7 \\
\hline
8 & Motivation \& Relevanz & 10\% & 9/10 & Alltagsrelevanz: Befinden erfragen \\
\hline
9 & Inklusion (UDL) & 10\% & 9/10 & Zeitzugabe, mündlich/schriftlich \\
\hline
10 & Praktikabilität & 5\% & 9/10 & 45+45 Min realistisch \\
\hline
\multicolumn{3}{|r|}{\textbf{GESAMT:}} & \textbf{9.1/10} & \textcolor{successgreen}{\textbf{FREIGABE}} \\
\hline
\end{tabular}
\end{scorebox}

\end{document}
